% STATUS: written 2026-11-04 (Bell and Daif 1994, Bresar and Miers 1994, Qi and Hou 2010, Bai and Du 2014 read in full)
% Sources read in full: Bell and Daif 1994 (pp. 443--447); Bresar and Miers 1994 (pp. 457--460); Qi and Hou 2010 (pp. 2790--2796); Bai and Du 2014 (pp. 733--742).
\section{Nonlinear and nonadditive Lie-type maps}\label{sec:lie-type:nonlinear}

The maps of Sections~\ref{sec:lie-type:martindale} to
\ref{sec:lie-type:triple-and-n} are additive or linear. The next step drops
even additivity. % claim-ok: section road-map, the plan of the book
Bell and Daif began this line with a definition and with two
commutativity theorems \cite[p.~443]{BellDaif1994}: a map $f\colon R\to R$ is
\emph{strong commutativity preserving} on a subset $S$ when
$[f(x),f(y)]=[x,y]$ for all $x,y\in S$ (Theorem~\ref{thm:lie:belldaif-derivation}).

\begin{theorem}[{\cite[Theorem~1, p.~444]{BellDaif1994}}]\label{thm:lie:belldaif-derivation}
Let $R$ be a semiprime ring and $U$ a nonzero right ideal. If $R$ admits a
derivation $D$ which is strong commutativity preserving on $U$, then $U$ is
contained in the center of $R$.
\end{theorem}

The second theorem of the pair keeps the map and replaces the derivation by
an endomorphism (Theorem~\ref{thm:lie:belldaif-endo}) \cite[p.~446]{BellDaif1994}.

\begin{theorem}[{\cite[Theorem~2, p.~446]{BellDaif1994}}]\label{thm:lie:belldaif-endo}
Let $R$ be a prime ring and $U$ an essential right ideal of $R$. If $R$
admits a nonidentity endomorphism $T$ which is strong commutativity
preserving on $U$, then $R$ is commutative.
\end{theorem}

Once such a map is allowed to be only additive on the whole ring, and
involves the extended centroid, the answers stop being commutativity
theorems and become form theorems. % claim-ok: section road-map, the plan of the book
Brešar and Miers gave the first one (Theorem~\ref{thm:lie:bresar-miers})
\cite[p.~457]{BresarMiers1994}.\begin{theorem}[{\cite[Theorem~1, p.~458]{BresarMiers1994}}]\label{thm:lie:bresar-miers}
Let $R$ be a semiprime ring with extended centroid $C$. Suppose that an
additive map $f\colon R\to R$ satisfies $[f(x),f(y)]=[x,y]$ for all
$x,y\in R$. Then $f$ is of the form $f(x)=\lambda x+\xi(x)$, where
$\lambda\in C$, $\lambda^{2}=1$, and $\xi$ is an additive map of $R$ into
$C$.
\end{theorem}

For prime rings with a nontrivial idempotent and surjective maps, no
additivity is needed: Qi and Hou obtained the same form with
$\lambda=\pm1$ (Theorem~\ref{thm:lie:qihou}) \cite[p.~2791]{QiHou2010}.

\begin{theorem}[{\cite[Theorem~2.1, p.~2791]{QiHou2010}}]\label{thm:lie:qihou}
Let $R$ be a unital prime ring containing a nontrivial idempotent. Assume
that $\varPhi\colon R\to R$ is a surjective map. Then $\varPhi$ is strong
commutativity preserving if and only if there exist a map $f\colon R\to Z(R)$
and an element $\lambda\in\{1,-1\}$ such that $\varPhi(A)=\lambda A+f(A)$
for all $A\in R$.
\end{theorem}

The factor von Neumann algebra case of their theorem is the scalar special
case, with $Z(R)$ replaced by the scalar operators \cite[p.~2795]{QiHou2010}.
Bai and Du then removed the prime hypothesis, at the price of asking the
idempotent to control the ring on both sides \cite[p.~733]{BaiDu2014}.

\begin{theorem}[{\cite[Theorem~1, p.~733]{BaiDu2014}}]\label{thm:lie:baidu-general}
Let $R$ be a unital ring having an idempotent element $e$ which satisfies
\begin{enumerate}
\item[(i)] $aRe=0$ implies $a=0$,
\item[(ii)] $aR(1-e)=0$ implies $a=0$.
\end{enumerate}
Assume that $f\colon R\to R$ is a surjective map. Then $f$ satisfies
$[f(x),f(y)]=[x,y]$ for all $x,y\in R$ if and only if $f$ is of the form
$f(x)=\alpha x+\xi(x)$ for all $x\in R$, where $\alpha\in Z(R)$,
$\alpha^{2}=1$, and $\xi$ is a map from $R$ into $Z(R)$.
\end{theorem}

The power of the idempotent hypothesis is visible in the operator case
\cite[p.~741]{BaiDu2014}. A von Neumann algebra with no central summands of type $I_{1}$ carries a
core-free projection, and that projection plays the role of the idempotent
of Theorem~\ref{thm:lie:baidu-general} \cite[p.~741]{BaiDu2014} (Corollary~\ref{cor:lie:baidu-vna}).

\begin{corollary}[{\cite[Corollary~2, pp.~734--735]{BaiDu2014}}]\label{cor:lie:baidu-vna}
Let $M$ be a von Neumann algebra with no central summands of type $I_{1}$.
Assume that $f\colon M\to M$ is a surjective map. Then $f$ satisfies
$[f(x),f(y)]=[x,y]$ for all $x,y\in M$ if and only if $f$ is of the form
$f(x)=\alpha x+\xi(x)$ for all $x\in M$, where $\alpha\in Z(M)$,
$\alpha^{2}=I$, and $\xi$ is a map from $M$ into $Z(M)$.
\end{corollary}

The pattern of the first four sections survives the loss of additivity:
again $f$ equals a scalar multiple, here $\alpha x$ with $\alpha^{2}=1$,
plus a central remainder $\xi(x)$ \cite[p.~733]{BaiDu2014}. The line runs
from a commutativity theorem to the same standard form that
Chapter~\ref{ch:jordan-type} and
Section~\ref{sec:lie-type:martindale} met for additive maps.

A $*$-algebra carries one more product of this kind, and the maps of this
section have an analogue for it \cite[p.~127]{AliAliTasleem2026}.

\begin{definition}[{\cite[p.~127]{AliAliTasleem2026}}]\label{def:lie:bi-skew-product}
Let $E$ be a $*$-algebra over the complex field $\mathbb{C}$. For $U,V\in E$
the \emph{left bi-skew Lie product} is $U\diamond V=U^{*}V-V^{*}U$. The product
of three elements is read as the left product applied twice, and the proof of
Lemma~11 records the identity
$U^{*}\diamond V\diamond\tfrac{I}{2}=-UV+V^{*}U^{*}$ for all $U,V\in E$
\cite[p.~136]{AliAliTasleem2026}.
\end{definition}

A family of maps that obeys the product rule of a higher derivation along this
product is a bi-skew Lie triple higher derivation
\cite[p.~128]{AliAliTasleem2026}. The sequence of
Definition~\ref{def:higher:hd} is the model for the name
\cite[p.~128]{AliAliTasleem2026}. Ali, Ali and Tasleem prove that no additivity
has to be assumed on a prime $*$-algebra that has a nontrivial projection
\cite[Main Theorem, p.~128]{AliAliTasleem2026}.

\begin{theorem}[{\cite[Main Theorem, p.~128]{AliAliTasleem2026}}]\label{thm:lie:ali-bi-skew-triple}
Let $E$ be a prime $*$-algebra with unity $I$ and a nontrivial projection $P$.
Let $\Delta=\{\zeta_{n}\}_{n\in\mathbb{N}}$ be a family of maps
$\zeta_{n}\colon E\to E$, not necessarily linear, with
$\zeta_{0}=\mathrm{id}_{E}$, such that
\[
\zeta_{n}(U\diamond V\diamond W)
=\sum_{p+q+r=n}\zeta_{p}(U)\diamond\zeta_{q}(V)\diamond\zeta_{r}(W)
\]
for all $U,V,W\in E$ and all $n\in\mathbb{N}$. Then $\Delta$ is an additive
$*$-higher derivation of $E$, provided $\zeta_{n}(\beta I/2)$ is self-adjoint
for $\beta\in\{1,i\}$ and every $n$.
\end{theorem}

The conclusion in full is the last lemma of the paper (Lemma~\ref{lem:lie:ali-bi-skew-additive}) \cite[p.~136]{AliAliTasleem2026}.

\begin{lemma}[{\cite[Lemma~11, p.~136]{AliAliTasleem2026}}]\label{lem:lie:ali-bi-skew-additive}
Under the hypotheses of Theorem~\ref{thm:lie:ali-bi-skew-triple}, each
$\zeta_{n}$ is additive, and
$\zeta_{n}(UV)=\sum_{p+q=n}\zeta_{p}(U)\zeta_{q}(V)$ and
$\zeta_{n}(U^{*})=\zeta_{n}(U)^{*}$ hold for all $U,V\in E$ and all $n\in\mathbb{N}$.
\end{lemma}

So a map that respects the triple bi-skew product at every order is additive and
multiplicative at every order \cite[p.~136]{AliAliTasleem2026}. The sequence of
maps and its own product rule are the subject of
Section~\ref{sec:higher-multi:hasse-schmidt}.

Motivated by Miers' work, nonlinear Lie-type derivations of a von Neumann algebra are characterised, which is the nonlinear form of the Lie-type question the chapter treats~\cite{WlNonlinearLieTypeDerivationsOfV2013}.

For a unital $*$-algebra containing two nontrivial projections, Ali, Haetinger, Naz and Tasleem show that every nonlinear mixed skew bi-skew Jordan $n$-derivation is an additive $*$-derivation, and extend the result to prime $*$-algebras, factor von Neumann algebras and standard operator algebras~\cite{WlNonlinearMixedSkewBiSkewJordan2026}.

Ashraf, Akhter and Ansari show for a complex $*$-algebra that a map $D$ satisfying the $n$-fold bi-skew Jordan derivation rule is a nonlinear bi-skew Jordan $n$-derivation exactly when it is an additive $*$-derivation, and prove the corresponding statement for the generalized notion they introduce~\cite{WlNonlinearGeneralizedBiSkewJord2024}.

For a unital $*$-algebra, Alam and Alshanquiti study nonlinear maps that satisfy the skew Lie-type rule built from the skew Jordan product $A\bullet B=AB+BA^*$ and the skew Lie product $[A,B]_*=AB-BA^*$, and characterise them~\cite{WlNonlinearMixedSkewLieTypeDeriv2026}.

Ansari, Alam, Alharbi and Mohamed study nonlinear mixed bi-skew Jordan and bi-skew Lie $n$-derivations on $*$-algebras and characterise them~\cite{WlANoteOnNonlinearMixedBiSkewJor2026}.

Rehman, Nisar and Nazim study nonlinear mixed Jordan triple derivations on $*$-algebras and show that they have the standard form on a unital $*$-algebra containing a nontrivial projection~\cite{WlANoteOnNonlinearMixedJordanTri2022}.
